\section{C2 Block-Dictionary Compaction}
\label{app:c2-block-dictionary}

This appendix records the C2 refinement implemented in the validation repository.  The original compaction layer publishes one linearly homomorphic code-LPN encryption of each coordinate of the extended GSW secret $\tilde{s}\in\F_q^N$.  Compacting a final row $r\in\F_q^N$ then evaluates
\[
  \sum_{i=1}^N r_i \cdot \Enc(\tilde{s}_i),
\]
so the number of accumulated compaction-noise terms is the row support $B=\|r\|_0$.  Since q-ary symmetric noise piles up as
\[
  \eta_B = \frac{q-1}{q}\left(1-\left(1-\frac{q\eta_c}{q-1}\right)^B\right),
\]
large $B$ forces a very small compaction noise rate $\eta_c$.

The C2 block-dictionary compactor trades public-key size for fewer accumulated compaction-noise terms.  Partition $\tilde{s}$ into blocks
\[
  \tilde{s}=(\tilde{s}^{(1)},\ldots,\tilde{s}^{(L)}),
  \qquad \tilde{s}^{(j)}\in\F_q^{\ell_j},
\]
where $\ell_j\leq \ell$ and $L=\lceil N/\ell\rceil$.  For every block $j$ and every nonzero coefficient vector $u\in\F_q^{\ell_j}$, publish
\[
  \mathsf{EK}^{\mathsf{C2}}_{j,u}=\CLPN\Enc_r(\langle u,\tilde{s}^{(j)}\rangle).
\]
Given a final GSW row $r$, write it in matching blocks $r=(r^{(1)},\ldots,r^{(L)})$ and output
\[
  \mathsf{Compact}_{\mathsf{C2}}(r)
  =\sum_{j:r^{(j)}\neq 0}\mathsf{EK}^{\mathsf{C2}}_{j,r^{(j)}}.
\]
By linear homomorphism of the CLPN layer, the plaintext inside the compacted ciphertext is
\[
  \sum_{j:r^{(j)}\neq 0}\langle r^{(j)},\tilde{s}^{(j)}\rangle
  =r\cdot\tilde{s}.
\]
Thus C2 preserves the compaction semantics.

The number of compaction-noise terms becomes
\[
  t_{\mathsf{C2}}(r)=\#\{j:r^{(j)}\neq 0\}
  \leq \min\{\|r\|_0,\lceil N/\ell\rceil\}.
\]
For dense rows this improves the compaction-noise term by about a factor of $\ell$.  For sparse rows there may be no improvement.  The public-key cost is
\[
  \sum_{j=1}^L(q^{\ell_j}-1)
\]
CLPN ciphertexts, so this construction is only plausible for small $q$ and small $\ell$, or for CRT-laned designs over small prime fields.  This is consistent with the code-based secure aggregation literature, where CRT-style decompositions are used to reduce communication in LPN-based homomorphic aggregation \cite{BitzerEggerLiuWachterZeh2026}.

The security proof remains a hybrid argument from CLPN semantic security: the evaluator sees additional CLPN encryptions of linear forms of $\tilde{s}$, but these messages are hidden under the compaction secret.  The important change is quantitative rather than syntactic: C2 increases the number of public CLPN samples, so the attack surface must be re-estimated.  In the validation repository the C2 estimator therefore reports both the reduced compaction-noise count and the enlarged public sample surface.

The executable validation reports that, for the toy setting $q=7$, $N=13$, and $\ell=2$, the block dictionary contains 294 ciphertexts and correctly evaluates clean multiplication and $xy+z$ examples.  For the larger design-screening setting $q=7$, $N=513$, and $\ell=3$, the dictionary contains 58,482 CLPN ciphertexts.  At multiplicative depth two, the worst-case compaction terms drop from 513 to 171, but the public sample surface becomes very large.  Therefore C2 should be treated as a research refinement, not as a certified parameter set.
